\begin{solapartat}
Les reaccions nuclears conserven la massa i la càrrega.
% A l'original es llegeix «c ==112» (doble signe d'igualtat), i així s'ha transcrit.
\[ \left.\begin{array}{l}
\left.\begin{array}{l} \text{massa inicial}: 208 + 70 = 278\\ \text{massa final}: 277 + d = 277 + d \end{array}\right\}d = 1\\[3ex]
\left.\begin{array}{l} \text{càrrega inicial}: a + b\\ \text{càrrega final}: c + e\\ a = 82, b = 30, c == 112 \end{array}\right\}e = 0
\end{array}\right\}
{}^{208}_{82}\mathit{Pb} + {}^{70}_{30}\mathit{Zn} \rightarrow {}^{277}_{112}\mathit{Cn} + {}^{1}_{0}n\ \text{(neutró)} \quad \text{\punts{0,4}} \]

\punts{0,6} El període de semidesintegració, $T_{1/2}$, és el temps que ha de passar per reduir-se a la meitat la quantitat d'una substància radioactiva, en aquest cas $T_{1/2} = 0,17\ \mathrm{ms} \Rightarrow$ En $t = 1$ minut no hi haurà res $\rightarrow 0\,\%$ de copernici.
\end{solapartat}

\begin{solapartat}
\punts{1}
\begin{minipage}[c]{0.45\linewidth}
\begin{align*}
&{}^{277}_{112}\mathit{Cn} \rightarrow {}^{273}_{110}\mathit{Ds} + {}^{4}_{2}\alpha\\
&{}^{273}_{110}\mathit{Ds} \rightarrow {}^{269}_{108}\mathit{Hs} + {}^{4}_{2}\alpha\\
&{}^{269}_{108}\mathit{Hs} \rightarrow {}^{265}_{106}\mathit{Sg} + {}^{4}_{2}\alpha\\
&{}^{265}_{106}\mathit{Sg} \rightarrow {}^{261}_{104}\mathit{Rf} + {}^{4}_{2}\alpha\\
&{}^{261}_{104}\mathit{Rf} \rightarrow {}^{257}_{102}\mathit{No} + {}^{4}_{2}\alpha\\
&{}^{257}_{102}\mathit{No} \rightarrow {}^{253}_{100}\mathit{Fm} + {}^{4}_{2}\alpha
\end{align*}
\end{minipage}\hfill
\begin{minipage}[c]{0.4\linewidth}
Es produeix una desintegració alfa repetida.
\end{minipage}

% A l'original «0,2 p.» i «zero» són en vermell.
Per cada desintegració incompleta o incorrecta es penalitzarà amb \punts{0,2}\unskip. Si totes les reaccions són incorrectes, la puntuació serà de \textcolor{red!70!black}{zero}.
\end{solapartat}
